% Options for packages loaded elsewhere
\PassOptionsToPackage{unicode}{hyperref}
\PassOptionsToPackage{hyphens}{url}
\documentclass[
]{article}
\usepackage{xcolor}
\usepackage{amsmath,amssymb}
\usepackage[margin=1in]{geometry}
\setcounter{secnumdepth}{4}
\usepackage{iftex}
\ifPDFTeX
  \usepackage[T1]{fontenc}
  \usepackage[utf8]{inputenc}
  \usepackage{textcomp} % provide euro and other symbols
\else % if luatex or xetex
  \usepackage{unicode-math} % this also loads fontspec
  \defaultfontfeatures{Scale=MatchLowercase}
  \defaultfontfeatures[\rmfamily]{Ligatures=TeX,Scale=1}
\fi
\usepackage{lmodern}
\ifPDFTeX\else
  % xetex/luatex font selection
\fi
% Use upquote if available, for straight quotes in verbatim environments
\IfFileExists{upquote.sty}{\usepackage{upquote}}{}
\IfFileExists{microtype.sty}{% use microtype if available
  \usepackage[]{microtype}
  \UseMicrotypeSet[protrusion]{basicmath} % disable protrusion for tt fonts
}{}
\makeatletter
\@ifundefined{KOMAClassName}{% if non-KOMA class
  \IfFileExists{parskip.sty}{%
    \usepackage{parskip}
  }{% else
    \setlength{\parindent}{0pt}
    \setlength{\parskip}{6pt plus 2pt minus 1pt}}
}{% if KOMA class
  \KOMAoptions{parskip=half}}
\makeatother
\usepackage{color}
\usepackage{fancyvrb}
\newcommand{\VerbBar}{|}
\newcommand{\VERB}{\Verb[commandchars=\\\{\}]}
\DefineVerbatimEnvironment{Highlighting}{Verbatim}{commandchars=\\\{\}}
% Add ',fontsize=\small' for more characters per line
\newenvironment{Shaded}{}{}
\newcommand{\AlertTok}[1]{\textcolor[rgb]{1.00,0.00,0.00}{\textbf{#1}}}
\newcommand{\AnnotationTok}[1]{\textcolor[rgb]{0.38,0.63,0.69}{\textbf{\textit{#1}}}}
\newcommand{\AttributeTok}[1]{\textcolor[rgb]{0.49,0.56,0.16}{#1}}
\newcommand{\BaseNTok}[1]{\textcolor[rgb]{0.25,0.63,0.44}{#1}}
\newcommand{\BuiltInTok}[1]{\textcolor[rgb]{0.00,0.50,0.00}{#1}}
\newcommand{\CharTok}[1]{\textcolor[rgb]{0.25,0.44,0.63}{#1}}
\newcommand{\CommentTok}[1]{\textcolor[rgb]{0.38,0.63,0.69}{\textit{#1}}}
\newcommand{\CommentVarTok}[1]{\textcolor[rgb]{0.38,0.63,0.69}{\textbf{\textit{#1}}}}
\newcommand{\ConstantTok}[1]{\textcolor[rgb]{0.53,0.00,0.00}{#1}}
\newcommand{\ControlFlowTok}[1]{\textcolor[rgb]{0.00,0.44,0.13}{\textbf{#1}}}
\newcommand{\DataTypeTok}[1]{\textcolor[rgb]{0.56,0.13,0.00}{#1}}
\newcommand{\DecValTok}[1]{\textcolor[rgb]{0.25,0.63,0.44}{#1}}
\newcommand{\DocumentationTok}[1]{\textcolor[rgb]{0.73,0.13,0.13}{\textit{#1}}}
\newcommand{\ErrorTok}[1]{\textcolor[rgb]{1.00,0.00,0.00}{\textbf{#1}}}
\newcommand{\ExtensionTok}[1]{#1}
\newcommand{\FloatTok}[1]{\textcolor[rgb]{0.25,0.63,0.44}{#1}}
\newcommand{\FunctionTok}[1]{\textcolor[rgb]{0.02,0.16,0.49}{#1}}
\newcommand{\ImportTok}[1]{\textcolor[rgb]{0.00,0.50,0.00}{\textbf{#1}}}
\newcommand{\InformationTok}[1]{\textcolor[rgb]{0.38,0.63,0.69}{\textbf{\textit{#1}}}}
\newcommand{\KeywordTok}[1]{\textcolor[rgb]{0.00,0.44,0.13}{\textbf{#1}}}
\newcommand{\NormalTok}[1]{#1}
\newcommand{\OperatorTok}[1]{\textcolor[rgb]{0.40,0.40,0.40}{#1}}
\newcommand{\OtherTok}[1]{\textcolor[rgb]{0.00,0.44,0.13}{#1}}
\newcommand{\PreprocessorTok}[1]{\textcolor[rgb]{0.74,0.48,0.00}{#1}}
\newcommand{\RegionMarkerTok}[1]{#1}
\newcommand{\SpecialCharTok}[1]{\textcolor[rgb]{0.25,0.44,0.63}{#1}}
\newcommand{\SpecialStringTok}[1]{\textcolor[rgb]{0.73,0.40,0.53}{#1}}
\newcommand{\StringTok}[1]{\textcolor[rgb]{0.25,0.44,0.63}{#1}}
\newcommand{\VariableTok}[1]{\textcolor[rgb]{0.10,0.09,0.49}{#1}}
\newcommand{\VerbatimStringTok}[1]{\textcolor[rgb]{0.25,0.44,0.63}{#1}}
\newcommand{\WarningTok}[1]{\textcolor[rgb]{0.38,0.63,0.69}{\textbf{\textit{#1}}}}
\usepackage{longtable,booktabs,array}
\newcounter{none} % for unnumbered tables
\usepackage{calc} % for calculating minipage widths
% Correct order of tables after \paragraph or \subparagraph
\usepackage{etoolbox}
\makeatletter
\patchcmd\longtable{\par}{\if@noskipsec\mbox{}\fi\par}{}{}
\makeatother
% Allow footnotes in longtable head/foot
\IfFileExists{footnotehyper.sty}{\usepackage{footnotehyper}}{\usepackage{footnote}}
\makesavenoteenv{longtable}
\setlength{\emergencystretch}{3em} % prevent overfull lines
\providecommand{\tightlist}{%
  \setlength{\itemsep}{0pt}\setlength{\parskip}{0pt}}
\usepackage{bookmark}
\IfFileExists{xurl.sty}{\usepackage{xurl}}{} % add URL line breaks if available
\urlstyle{same}
\hypersetup{
  hidelinks,
  pdfcreator={LaTeX via pandoc}}
\newcommand{\indextopic}[2]{\noindent\hyperref[#1]{\textbf{#2}}\par}
\newcommand{\indexentry}[2]{\noindent\hspace*{1.5em}\hyperref[#1]{#2}\nobreak\hfill\hyperref[#1]{\pageref*{#1}}\par}

\author{}
\date{}

\begin{document}
\section*{Sentiment Analysis Documentation}
\indextopic{topic-overview}{1. Overview}
\vspace{0.25em}
\indexentry{file-project-overview}{1.1 Project Overview}
\indexentry{file-configuration-cli}{1.2 Configuration \& CLI}
\vspace{0.9em}

\indextopic{topic-pipelines}{2. Pipelines}
\vspace{0.25em}
\indexentry{file-article-pipeline}{2.1 Article Pipeline}
\indexentry{file-youtube-pipeline}{2.2 YouTube Pipeline}
\vspace{0.9em}

\indextopic{topic-database}{3. Database}
\vspace{0.25em}
\indexentry{file-articles-schema}{3.1 Articles Schema}
\indexentry{file-youtube-schema}{3.2 YouTube Schema}
\vspace{0.9em}

\indextopic{topic-llm-usage}{4. LLM Usage}
\vspace{0.25em}
\indexentry{file-gemini-kg}{4.1 Gemini --- Knowledge Graph Extraction}
\indexentry{file-marianmt-translation}{4.2 MarianMT --- Bangla Translation}
\newpage

\phantomsection\label{topic-overview}
\section{Overview}
\phantomsection\label{file-project-overview}
\subsection{Project Overview}\label{sentiment-analysis--project-overview}

A unified data ingestion pipeline that scrapes \textbf{web news
articles} and \textbf{YouTube channels} into a shared PostgreSQL
database. Originally two separate services, now merged into a single
monorepo with shared configuration, a shared database layer, and a
unified CLI.

\begin{center}\rule{0.5\linewidth}{0.5pt}\end{center}

\subsubsection{What It Does}\label{what-it-does}

\textbf{Article Pipeline} Takes a news article URL, scrapes the full
page using a headless browser, detects the language (with automatic
Bangla → English translation), and sends the body text to Google Gemini
to extract a structured knowledge graph of named entities and their
relationships. Everything --- the article, its translation, and the
knowledge graph --- is stored in the database.

\textbf{YouTube Pipeline} Takes a registered YouTube channel, enumerates
all its videos via yt-dlp, and processes each one as a Celery task. For
every video it downloads the file and captions to NAS, fetches comments
via the YouTube Data API, and extracts linked external articles and
related YouTube videos from descriptions, pinned comments, end-screens,
and sidebar suggestions.

\begin{center}\rule{0.5\linewidth}{0.5pt}\end{center}

\subsubsection{Tech Stack}\label{tech-stack}

{\def\LTcaptype{none} % do not increment counter
\begin{longtable}[]{@{}ll@{}}
\toprule\noalign{}
Layer & Technology \\
\midrule\noalign{}
\endhead
\bottomrule\noalign{}
\endlastfoot
Scraping & Playwright (articles), yt-dlp (YouTube) \\
NLP & Google Gemini (KG extraction), Helsinki-NLP MarianMT (Bangla
translation) \\
Task Queue & Celery + Redis \\
Database & PostgreSQL + SQLAlchemy 2.0 + Alembic \\
Language & Python 3.12, asyncio \\
\end{longtable}
}

\begin{center}\rule{0.5\linewidth}{0.5pt}\end{center}

\subsubsection{Shared Infrastructure}\label{shared-infrastructure}

Both pipelines share:

\begin{itemize}
\tightlist
\item
  A single PostgreSQL database with two service-owned schemas ---
  \texttt{articles} and \texttt{youtube}
\item
  \texttt{sentiment\_analysis/database.py} --- one SQLAlchemy engine and
  \texttt{get\_session()} context manager used everywhere
\item
  \texttt{sentiment\_analysis/config.py} --- single source of truth for
  all environment variables, loaded from \texttt{.env}
\item
  One Alembic migration chain covering both schemas
\end{itemize}

\begin{center}\rule{0.5\linewidth}{0.5pt}\end{center}

\subsubsection{Project Structure}\label{project-structure}

\begin{verbatim}
sentiment_analysis/
├── config.py          — all env vars + logging setup
├── database.py        — shared SQLAlchemy engine and get_session()
├── db_migrations.py   — Alembic runner, auto-creates DB if missing
├── cli.py             — argument parsing and command dispatch
├── articles/          — article ingestion sub-package
│   ├── scraper.py     — Playwright + Trafilatura + BeautifulSoup
│   ├── translator.py  — Bangla → English (MarianMT)
│   ├── extractor.py   — Gemini KG extraction
│   ├── db.py          — persistence layer
│   ├── models.py      — ORM models (articles schema)
│   └── service.py     — ingest_url() orchestrator
└── youtube/           — YouTube ingestion sub-package
    ├── tasks.py        — Celery tasks (scrape_channel, process_video)
    ├── celery_app.py   — Celery + Redis config
    ├── service.py      — CLI command handlers
    ├── models.py       — ORM models (youtube schema)
    └── scraper/
        ├── ytdlp.py    — enumerate, metadata, captions, download
        ├── comments.py — YouTube Data API comments
        ├── innertube.py — suggested + end-screen video IDs
        └── articles.py — link extraction from descriptions
alembic/                — single migration chain for both schemas
main.py                 — CLI entrypoint
\end{verbatim}

\begin{center}\rule{0.5\linewidth}{0.5pt}\end{center}

\subsubsection{Quick Start}\label{quick-start}

\begin{Shaded}
\begin{Highlighting}[]
\FunctionTok{cp}\NormalTok{ .env.example .env          }\CommentTok{\# fill in credentials}
\ExtensionTok{python}\NormalTok{ main.py init           }\CommentTok{\# create DB + run all migrations}

\CommentTok{\# Ingest an article}
\ExtensionTok{python}\NormalTok{ main.py articles ingest }\StringTok{"https://example.com/article"}

\CommentTok{\# Register and scrape a YouTube channel}
\ExtensionTok{python}\NormalTok{ main.py youtube add{-}channel }\AttributeTok{{-}{-}channel{-}id}\NormalTok{ UCxxxx }\AttributeTok{{-}{-}name} \StringTok{"Name"} \AttributeTok{{-}{-}url} \StringTok{"https://..."}
\ExtensionTok{celery} \AttributeTok{{-}A}\NormalTok{ sentiment\_analysis.youtube.celery\_app worker }\AttributeTok{{-}{-}loglevel}\OperatorTok{=}\NormalTok{info }\AttributeTok{{-}{-}concurrency}\OperatorTok{=}\NormalTok{4}
\ExtensionTok{python}\NormalTok{ main.py youtube scrape }\AttributeTok{{-}{-}channel{-}id}\NormalTok{ UCxxxx}
\end{Highlighting}
\end{Shaded}

\phantomsection\label{file-configuration-cli}
\subsection{Configuration \& CLI}\label{configuration--cli}

All configuration lives in \texttt{sentiment\_analysis/config.py} and is
loaded from a \texttt{.env} file at startup via \texttt{python-dotenv}.
Copy \texttt{.env.example} to \texttt{.env} and fill in the required
values before running anything.

\begin{center}\rule{0.5\linewidth}{0.5pt}\end{center}

\subsubsection{Environment Variables}\label{environment-variables}

{\def\LTcaptype{none} % do not increment counter
\small
\setlength{\tabcolsep}{4pt}
\begin{longtable}[]{@{}lll@{}}
\toprule\noalign{}
Variable & Default & Description \\
\midrule\noalign{}
\endhead
\bottomrule\noalign{}
\endlastfoot
\texttt{DATABASE\_URL} & --- & PostgreSQL connection string ---
required \\
\texttt{REDIS\_URL} & \texttt{redis://localhost:6379/0} & Celery broker
and result backend \\
\texttt{YOUTUBE\_API\_KEY} & --- & YouTube Data API v3 key --- required
for comment fetching \\
\texttt{GEMINI\_API\_KEY} & --- & Google Gemini API key --- required for
article KG extraction \\
\texttt{GEMINI\_MODEL} & \texttt{gemini-3.1-pro-preview} & Primary model
used for KG extraction \\
\texttt{GEMINI\_FALLBACK\_MODEL} & \texttt{gemini-2.5-pro} & Fallback
model if primary fails \\
\texttt{GEMINI\_MAX\_TOKENS} & \texttt{8000} & Max tokens of body text
sent to Gemini; longer text is truncated \\
\texttt{TRANSLATION\_MODEL} & \texttt{Helsinki-NLP/opus-mt-bn-en} &
HuggingFace model for Bangla → English translation \\
\texttt{NAS\_BASE\_PATH} & \texttt{/Volumes/Downloaded\ YT\ Videos} &
Mount point for downloaded video files \\
\texttt{CAPTIONS\_BASE\_PATH} & \texttt{/Volumes/Captions-\ YT\ Videos}
& Mount point for caption (VTT) files \\
\texttt{REQUEST\_DELAY\_SECONDS} & \texttt{3} & Delay between yt-dlp and
InnerTube API calls to avoid rate limits \\
\texttt{MAX\_RETRIES} & \texttt{3} & Max Celery task retry attempts
before giving up \\
\texttt{VIDEO\_QUALITY} & \texttt{1080} & Max video resolution in pixels
for yt-dlp downloads \\
\texttt{LOG\_LEVEL} & \texttt{INFO} & Python logging level
(\texttt{DEBUG}, \texttt{INFO}, \texttt{WARNING}, etc.) \\
\end{longtable}
}

\begin{center}\rule{0.5\linewidth}{0.5pt}\end{center}

\subsubsection{CLI Reference}\label{cli-reference}

The unified entrypoint is \texttt{python\ main.py}. Commands are
dispatched via \texttt{sentiment\_analysis/cli.py}.

\paragraph{Initialisation}\label{initialisation}

\begin{Shaded}
\begin{Highlighting}[]
\ExtensionTok{python}\NormalTok{ main.py init}
\end{Highlighting}
\end{Shaded}

Creates the database if it doesn\textquotesingle t exist, then runs all
Alembic migrations. Run this once before using either pipeline.

\begin{center}\rule{0.5\linewidth}{0.5pt}\end{center}

\paragraph{Articles}\label{articles}

\begin{Shaded}
\begin{Highlighting}[]
\ExtensionTok{python}\NormalTok{ main.py articles ingest }\StringTok{"\textless{}url\textgreater{}"}
\end{Highlighting}
\end{Shaded}

Runs the full article pipeline for the given URL --- scrape, translate
if needed, extract KG, save to DB. Prints the article hash on
completion. Exits early without error if the URL was already ingested.

\begin{center}\rule{0.5\linewidth}{0.5pt}\end{center}

\paragraph{YouTube}\label{youtube}

\begin{Shaded}
\begin{Highlighting}[]
\CommentTok{\# Register a channel (required before scraping)}
\ExtensionTok{python}\NormalTok{ main.py youtube add{-}channel }\DataTypeTok{\textbackslash{}}
  \AttributeTok{{-}{-}channel{-}id}\NormalTok{ UCxxxx }\DataTypeTok{\textbackslash{}}
  \AttributeTok{{-}{-}name} \StringTok{"Channel Name"} \DataTypeTok{\textbackslash{}}
  \AttributeTok{{-}{-}url} \StringTok{"https://www.youtube.com/@handle"}

\CommentTok{\# Queue a full channel scrape via Celery}
\ExtensionTok{python}\NormalTok{ main.py youtube scrape }\AttributeTok{{-}{-}channel{-}id}\NormalTok{ UCxxxx}

\CommentTok{\# Limit to N most recent videos (useful for testing)}
\ExtensionTok{python}\NormalTok{ main.py youtube scrape }\AttributeTok{{-}{-}channel{-}id}\NormalTok{ UCxxxx }\AttributeTok{{-}{-}limit}\NormalTok{ 50}

\CommentTok{\# Show processing status across all registered channels}
\ExtensionTok{python}\NormalTok{ main.py youtube status}

\CommentTok{\# Re{-}queue videos stuck in pending or processing}
\ExtensionTok{python}\NormalTok{ main.py youtube retry }\AttributeTok{{-}{-}channel{-}id}\NormalTok{ UCxxxx}
\end{Highlighting}
\end{Shaded}

\texttt{scrape} only dispatches tasks --- actual processing requires a
running Celery worker:

\begin{Shaded}
\begin{Highlighting}[]
\ExtensionTok{celery} \AttributeTok{{-}A}\NormalTok{ sentiment\_analysis.youtube.celery\_app worker }\AttributeTok{{-}{-}loglevel}\OperatorTok{=}\NormalTok{info }\AttributeTok{{-}{-}concurrency}\OperatorTok{=}\NormalTok{4}
\end{Highlighting}
\end{Shaded}

\phantomsection\label{file-article-pipeline}
\section{Article Pipeline}\label{article-pipeline}

\textbf{Entrypoint:}
\texttt{python\ main.py\ articles\ ingest\ \textless{}url\textgreater{}}
\textbf{Orchestrator:}
\texttt{sentiment\_analysis/articles/service.py\ →\ ingest\_url()}
\textbf{Execution:} fully async via \texttt{asyncio} --- required
because both Playwright and the Gemini client expose async APIs

\begin{center}\rule{0.5\linewidth}{0.5pt}\end{center}

\subsection{Flow}\label{flow}

\begin{verbatim}
URL → dedup check → scrape → detect language → translate (if Bangla)
    → save article → KG extraction (Gemini) → save KG
\end{verbatim}

\begin{center}\rule{0.5\linewidth}{0.5pt}\end{center}

\subsection{Steps}\label{steps}

\subsubsection{1. Deduplication}\label{1-deduplication}

Before anything runs, the pipeline queries \texttt{articles.articles}
for the given \texttt{source\_url}. If a match is found, the existing
hash is returned and the pipeline exits immediately --- no scraping, no
API calls.

\subsubsection{2. Scraping}\label{2-scraping}

\texttt{scraper.py} launches a headless Chromium browser via Playwright
with a realistic user-agent to avoid bot blocks. It waits for the page
to load, captures the full HTML, then:

\begin{itemize}
\tightlist
\item
  Passes the HTML to \textbf{Trafilatura} to extract clean body text
  (strips ads, nav, boilerplate)
\item
  Parses metadata with \textbf{BeautifulSoup} --- headline, author
  names, publish date, and images --- using a priority chain (JSON-LD →
  Open Graph meta → HTML fallbacks)
\end{itemize}

The scraper retries up to \textbf{3 times with exponential backoff} on
failure.

\subsubsection{3. Language Detection}\label{3-language-detection}

\texttt{langdetect} identifies the language of the body text. Only
Bangla (\texttt{bn}) triggers translation --- all other languages pass
through as-is.

\subsubsection{4. Translation (Bangla
only)}\label{4-translation-bangla-only}

If Bangla is detected, \texttt{translator.py} runs the headline, body,
and reporter names through \textbf{Helsinki-NLP/opus-mt-bn-en} (a
MarianMT model loaded lazily on first use). The model has a 512-token
limit, so text is split into paragraph and sentence-level chunks before
translating. Original Bangla text is preserved in \texttt{*\_original}
columns in the DB.

\subsubsection{5. Save Article}\label{5-save-article}

\texttt{db.py} generates a deterministic hash in the format
\texttt{\{source\}-\{ddmmyyyy\}-\{hhmmss\}-\{seq:03d\}} (e.g.
\texttt{bbc-08042026-143022-001}). The sequence number is determined
with \texttt{SELECT\ ...\ FOR\ UPDATE} to prevent collisions under
concurrent ingestion. The article row is then inserted.

\subsubsection{6. Knowledge Graph
Extraction}\label{6-knowledge-graph-extraction}

\texttt{extractor.py} sends the body text to \textbf{Google Gemini} with
a structured prompt instructing it to extract:

\begin{itemize}
\tightlist
\item
  \textbf{Entities (nodes):} PERSON, ORGANIZATION, LOCATION, COUNTRY,
  COMPANY, EVENT --- canonical names only
\item
  \textbf{Relations (edges):} directed, uppercase snake\_case labels
  like \texttt{VISITED}, \texttt{APPOINTED\_BY}, \texttt{ACQUIRED}
\end{itemize}

Text is truncated to \texttt{GEMINI\_MAX\_TOKENS} (default 8000) using
tiktoken before sending. If the primary model fails, it automatically
retries with the fallback model.

\subsubsection{7. Save Knowledge Graph}\label{7-save-knowledge-graph}

Elements and relations are upserted into \texttt{kg\_elements} and
\texttt{kg\_relations} using raw \texttt{ON\ CONFLICT} SQL. Any relation
referencing an entity not in the extracted elements list is silently
skipped (guards against Gemini hallucinations).

\begin{center}\rule{0.5\linewidth}{0.5pt}\end{center}

\subsection{Key Files}\label{key-files}

{\def\LTcaptype{none} % do not increment counter
\begin{longtable}[]{@{}ll@{}}
\toprule\noalign{}
File & Role \\
\midrule\noalign{}
\endhead
\bottomrule\noalign{}
\endlastfoot
\texttt{service.py} & Orchestrates the full pipeline end-to-end \\
\texttt{scraper.py} & Playwright fetch, Trafilatura extraction,
BeautifulSoup metadata parsing \\
\texttt{translator.py} & Lazy-loaded MarianMT model, chunk splitting,
Bangla → English \\
\texttt{extractor.py} & Gemini API call, primary/fallback model logic,
token truncation \\
\texttt{db.py} & Article insert with hash generation, KG upsert \\
\texttt{models.py} & \texttt{Article}, \texttt{KGElement},
\texttt{KGRelation} SQLAlchemy models \\
\end{longtable}
}

\phantomsection\label{file-youtube-pipeline}
\section{YouTube Pipeline}\label{youtube-pipeline}

\textbf{Entrypoint:}
\texttt{python\ main.py\ youtube\ scrape\ -\/-channel-id\ UCxxxx}
\textbf{Execution:} Celery distributed task queue backed by Redis
\textbf{Worker:}
\texttt{celery\ -A\ sentiment\_analysis.youtube.celery\_app\ worker\ -\/-loglevel=info\ -\/-concurrency=4}

\begin{center}\rule{0.5\linewidth}{0.5pt}\end{center}

\subsection{Flow}\label{flow-1}

\begin{verbatim}
CLI → scrape_channel task (Phase A)
          └─► process_video task × N (Phases B → E)
                  ├── B: metadata + captions
                  ├── C: video download → NAS
                  ├── D: comments (YouTube Data API)
                  └── E: linked articles + related videos
\end{verbatim}

Each video is an independent Celery task, so a channel with 500 videos
spawns 500 parallel tasks processed across however many workers are
running.

\begin{center}\rule{0.5\linewidth}{0.5pt}\end{center}

\subsection{Phases}\label{phases}

\subsubsection{\texorpdfstring{Phase A --- Channel Enumeration
(\texttt{tasks.py},
\texttt{ytdlp.py})}{Phase A --- Channel Enumeration (tasks.py, ytdlp.py)}}\label{phase-a--channel-enumeration-taskspy-ytdlppy}

yt-dlp lists all video IDs from the channel\textquotesingle s
\texttt{/videos} tab in flat-extract mode (fast --- no metadata fetched
yet). The list is reversed to chronological order, filtered against
video IDs already in the DB, and then:

\begin{enumerate}
\def\labelenumi{\arabic{enumi}.}
\tightlist
\item
  \texttt{pending} records are inserted into the DB for every new video
  \textbf{before} any Celery tasks are dispatched --- this ensures no
  video is silently lost if a task message is dropped
\item
  One \texttt{process\_video} task is dispatched per new video
\item
  \texttt{last\_scraped\_at} is updated on the channel record
\end{enumerate}

\subsubsection{\texorpdfstring{Phase B --- Metadata \& Captions
(\texttt{ytdlp.py})}{Phase B --- Metadata \& Captions (ytdlp.py)}}\label{phase-b--metadata--captions-ytdlppy}

yt-dlp fetches the full video metadata without downloading (title,
description, duration, view count, like count, thumbnail, upload date
and timestamp). The \texttt{internal\_id} is built from this:
\texttt{\{channel\}-\{ddmmyyyy\}-\{hhmmss\}-\{seq\}}.

Captions are downloaded separately in VTT format to
\texttt{\{CAPTIONS\_BASE\_PATH\}/\{channel\_id\}/}. Manual subtitles are
tried first; auto-generated captions are the fallback. The relative path
is stored in the DB as \texttt{\{"en":\ "path/to/file.en.vtt"\}}.

\subsubsection{\texorpdfstring{Phase C --- Video Download
(\texttt{ytdlp.py})}{Phase C --- Video Download (ytdlp.py)}}\label{phase-c--video-download-ytdlppy}

The video file is downloaded to
\texttt{\{NAS\_BASE\_PATH\}/\{channel\_id\}/} at up to
\texttt{VIDEO\_QUALITY} resolution (default 1080p), merged into MP4. The
path stored in the DB is relative to \texttt{NAS\_BASE\_PATH} so it
stays valid regardless of where the volume is mounted.

\subsubsection{\texorpdfstring{Phase D --- Comments
(\texttt{comments.py})}{Phase D --- Comments (comments.py)}}\label{phase-d--comments-commentspy}

The YouTube Data API v3 \texttt{commentThreads} endpoint is paginated
(100 per page) to collect all top-level comments and their replies. Each
comment stores the author, text, timestamp, and a \texttt{reply\_to}
field (null for top-level, parent comment ID for replies). Skips
gracefully if comments are disabled or \texttt{YOUTUBE\_API\_KEY} is
absent.

\subsubsection{\texorpdfstring{Phase E --- Linked Content Discovery
(\texttt{articles.py},
\texttt{innertube.py})}{Phase E --- Linked Content Discovery (articles.py, innertube.py)}}\label{phase-e--linked-content-discovery-articlespy-innertubepy}

Three types of related content are extracted and stored:

\begin{itemize}
\tightlist
\item
  \textbf{Linked articles} --- non-YouTube URLs from the description and
  pinned comment. Each URL is followed through redirects to resolve
  shorteners (e.g. \texttt{bbc.in} → full URL). Domain and page title
  are recorded.
\item
  \textbf{Related videos from description} --- YouTube video URLs found
  in the description text, stored as \texttt{description-linked}
\item
  \textbf{Suggested videos} --- sidebar suggestions fetched from
  YouTube\textquotesingle s internal InnerTube \texttt{/next} API,
  stored as \texttt{suggested}
\item
  \textbf{End-screen videos} --- video IDs extracted from
  \texttt{ytInitialPlayerResponse} in the watch page HTML, stored as
  \texttt{end-screen}
\end{itemize}

All sources are deduplicated before writing to the DB.

\begin{center}\rule{0.5\linewidth}{0.5pt}\end{center}

\subsection{Error Handling}\label{error-handling}

Non-fatal steps (captions, comments, Phase E) are individually wrapped
--- a failure in one logs a warning and the pipeline continues. The
video is still marked \texttt{processed}.

For \texttt{DownloadError} from yt-dlp: rate limits, SSL errors, and
ffmpeg failures retry with a 120-second delay. Permanently unavailable
videos (private, deleted) are logged and discarded cleanly.

All other errors trigger automatic Celery retries up to
\texttt{MAX\_RETRIES} times (default 3). \texttt{task\_acks\_late=True}
and \texttt{task\_reject\_on\_worker\_lost=True} ensure tasks are not
lost if a worker is killed mid-execution.

Stuck \texttt{pending} or \texttt{processing} videos can be re-queued
with:

\begin{Shaded}
\begin{Highlighting}[]
\ExtensionTok{python}\NormalTok{ main.py youtube retry }\AttributeTok{{-}{-}channel{-}id}\NormalTok{ UCxxxx}
\end{Highlighting}
\end{Shaded}

\begin{center}\rule{0.5\linewidth}{0.5pt}\end{center}

\subsection{Key Files}\label{key-files-1}

{\def\LTcaptype{none} % do not increment counter
\begin{longtable}[]{@{}ll@{}}
\toprule\noalign{}
File & Role \\
\midrule\noalign{}
\endhead
\bottomrule\noalign{}
\endlastfoot
\texttt{tasks.py} & \texttt{scrape\_channel} (Phase A) and
\texttt{process\_video} (B--E) Celery tasks \\
\texttt{celery\_app.py} & Celery app wired to Redis with reliability
settings \\
\texttt{service.py} & CLI handlers --- add-channel, scrape, status,
retry \\
\texttt{scraper/ytdlp.py} & Channel enumeration, metadata, captions,
video download \\
\texttt{scraper/comments.py} & YouTube Data API v3 comment pagination \\
\texttt{scraper/innertube.py} & Suggested and end-screen video IDs \\
\texttt{scraper/articles.py} & URL extraction and redirect resolution
from description/comments \\
\texttt{models.py} & \texttt{Channel}, \texttt{Video}, \texttt{Comment},
\texttt{RelatedVideo}, \texttt{LinkedArticle} \\
\end{longtable}
}

\phantomsection\label{file-articles-schema}
\section{Articles Schema}\label{articles-schema}

The \texttt{articles} schema stores scraped news article content
alongside the knowledge graph extracted from each article by Gemini. All
tables cascade-delete on article removal.

\begin{center}\rule{0.5\linewidth}{0.5pt}\end{center}

\subsection{\texorpdfstring{\texttt{articles.articles}}{articles.articles}}\label{articlesarticles}

One row per scraped news article. For Bangla sources, translated English
content and the original Bangla text are stored side by side.

{\def\LTcaptype{none} % do not increment counter
\begin{longtable}[]{@{}lll@{}}
\toprule\noalign{}
Column & Type & Notes \\
\midrule\noalign{}
\endhead
\bottomrule\noalign{}
\endlastfoot
\texttt{hash} & String & Unique ID:
\texttt{\{source\}-\{ddmmyyyy\}-\{hhmmss\}-\{seq:03d\}} \\
\texttt{source\_url} & Text & Unique --- used for deduplication \\
\texttt{headline} & Text & English (translated if source is Bangla) \\
\texttt{reporters} & Text{[}{]} & Array of author names \\
\texttt{body\_text} & Text & Cleaned article body \\
\texttt{media\_urls} & JSONB &
\texttt{{[}\{"url":\ "...",\ "type":\ "image"\}{]}} \\
\texttt{language} & String & Detected language code (e.g. \texttt{en},
\texttt{bn}) \\
\texttt{headline\_original} & Text & Bangla original headline --- null
if not translated \\
\texttt{body\_text\_original} & Text & Bangla original body --- null if
not translated \\
\texttt{reporters\_original} & Text{[}{]} & Bangla original reporter
names \\
\end{longtable}
}

Has a GIN full-text search index on
\texttt{headline\ \textbar{}\textbar{}\ body\_text}.

\begin{center}\rule{0.5\linewidth}{0.5pt}\end{center}

\subsection{\texorpdfstring{\texttt{articles.kg\_elements}}{articles.kg\_elements}}\label{articleskg_elements}

Named entity nodes extracted by Gemini. Each row is one distinct entity
within one article.

{\def\LTcaptype{none} % do not increment counter
\begin{longtable}[]{@{}ll@{}}
\toprule\noalign{}
Column & Notes \\
\midrule\noalign{}
\endhead
\bottomrule\noalign{}
\endlastfoot
\texttt{article\_hash} & FK → \texttt{articles.hash} \\
\texttt{name} & Canonical entity name (e.g. \texttt{"United\ Nations"},
not \texttt{"the\ UN"}) \\
\texttt{entity\_type} & One of:
\texttt{PERSON\ /\ ORGANIZATION\ /\ LOCATION\ /\ COUNTRY\ /\ COMPANY\ /\ EVENT} \\
\end{longtable}
}

Unique on \texttt{(article\_hash,\ name,\ entity\_type)} --- no
duplicate entities per article.

\begin{center}\rule{0.5\linewidth}{0.5pt}\end{center}

\subsection{\texorpdfstring{\texttt{articles.kg\_relations}}{articles.kg\_relations}}\label{articleskg_relations}

Directed edges between knowledge graph nodes within the same article.

{\def\LTcaptype{none} % do not increment counter
\begin{longtable}[]{@{}ll@{}}
\toprule\noalign{}
Column & Notes \\
\midrule\noalign{}
\endhead
\bottomrule\noalign{}
\endlastfoot
\texttt{article\_hash} & FK → \texttt{articles.hash} \\
\texttt{subject\_id} & FK → \texttt{kg\_elements.id} --- the "from"
node \\
\texttt{object\_id} & FK → \texttt{kg\_elements.id} --- the "to" node \\
\texttt{relation} & Uppercase snake\_case label (e.g.
\texttt{APPOINTED\_BY}, \texttt{HEADQUARTERED\_IN}) \\
\end{longtable}
}

Unique on
\texttt{(article\_hash,\ subject\_id,\ object\_id,\ relation)}.

\phantomsection\label{file-youtube-schema}
\section{YouTube Schema}\label{youtube-schema}

The \texttt{youtube} schema stores everything collected from YouTube
channels --- channel records, per-video metadata and file paths,
comments, related videos, and linked external articles. All tables
cascade-delete from \texttt{channels} down.

\begin{center}\rule{0.5\linewidth}{0.5pt}\end{center}

\subsection{\texorpdfstring{\texttt{youtube.channels}}{youtube.channels}}\label{youtubechannels}

Registered YouTube channels. A channel must be added here before it can
be scraped.

{\def\LTcaptype{none} % do not increment counter
\begin{longtable}[]{@{}ll@{}}
\toprule\noalign{}
Column & Notes \\
\midrule\noalign{}
\endhead
\bottomrule\noalign{}
\endlastfoot
\texttt{channel\_id} & YouTube\textquotesingle s channel ID string (e.g.
\texttt{UCxxxxxx}) --- unique \\
\texttt{name} & Human-readable display name \\
\texttt{url} & Channel URL (e.g.
\texttt{https://www.youtube.com/@bbcnews}) \\
\texttt{last\_scraped\_at} & UTC timestamp of the last successful
\texttt{scrape\_channel} run \\
\end{longtable}
}

\begin{center}\rule{0.5\linewidth}{0.5pt}\end{center}

\subsection{\texorpdfstring{\texttt{youtube.videos}}{youtube.videos}}\label{youtubevideos}

One row per video. Created as \texttt{pending} in Phase A and updated to
\texttt{processed} after all phases complete.

{\def\LTcaptype{none} % do not increment counter
\begin{longtable}[]{@{}ll@{}}
\toprule\noalign{}
Column & Notes \\
\midrule\noalign{}
\endhead
\bottomrule\noalign{}
\endlastfoot
\texttt{internal\_id} & Unique slug:
\texttt{\{channel\}-\{ddmmyyyy\}-\{hhmmss\}-\{seq\}} \\
\texttt{yt\_video\_id} & YouTube\textquotesingle s video ID ---
unique \\
\texttt{channel\_id} & FK → \texttt{channels.id} \\
\texttt{title}, \texttt{description}, \texttt{duration},
\texttt{view\_count}, \texttt{like\_count}, \texttt{thumbnail\_url} &
Metadata from yt-dlp \\
\texttt{nas\_file\_path} & Relative path to the downloaded video on
NAS \\
\texttt{caption\_nas\_path} & JSONB:
\texttt{\{"en":\ "UCxxxxxx/videoId.en.vtt"\}} \\
\texttt{status} & \texttt{pending\ →\ processing\ →\ processed} \\
\end{longtable}
}

\begin{center}\rule{0.5\linewidth}{0.5pt}\end{center}

\subsection{\texorpdfstring{\texttt{youtube.comments}}{youtube.comments}}\label{youtubecomments}

All comments and replies fetched from the YouTube Data API. Top-level
comments and replies share the same table, distinguished by
\texttt{reply\_to}.

{\def\LTcaptype{none} % do not increment counter
\begin{longtable}[]{@{}ll@{}}
\toprule\noalign{}
Column & Notes \\
\midrule\noalign{}
\endhead
\bottomrule\noalign{}
\endlastfoot
\texttt{video\_id} & FK → \texttt{videos.id} \\
\texttt{comment\_id} & YouTube\textquotesingle s comment ID ---
unique \\
\texttt{author}, \texttt{text}, \texttt{published\_at} & Comment
content \\
\texttt{reply\_to} & Parent comment ID --- null for top-level, set for
replies \\
\end{longtable}
}

\begin{center}\rule{0.5\linewidth}{0.5pt}\end{center}

\subsection{\texorpdfstring{\texttt{youtube.related\_videos}}{youtube.related\_videos}}\label{youtuberelated_videos}

YouTube videos linked to a given video, sourced from descriptions,
sidebar suggestions, and end-screens.

{\def\LTcaptype{none} % do not increment counter
\begin{longtable}[]{@{}ll@{}}
\toprule\noalign{}
Column & Notes \\
\midrule\noalign{}
\endhead
\bottomrule\noalign{}
\endlastfoot
\texttt{video\_id} & FK → \texttt{videos.id} \\
\texttt{related\_video\_id} & YouTube video ID of the related video \\
\texttt{url} & Full \texttt{https://www.youtube.com/watch?v=...} URL \\
\texttt{relation\_type} &
\texttt{suggested\ /\ end-screen\ /\ description-linked} \\
\end{longtable}
}

Unique on \texttt{(video\_id,\ related\_video\_id)}.

\begin{center}\rule{0.5\linewidth}{0.5pt}\end{center}

\subsection{\texorpdfstring{\texttt{youtube.linked\_articles}}{youtube.linked\_articles}}\label{youtubelinked_articles}

External (non-YouTube) article URLs found in a video\textquotesingle s
description or pinned comment. Redirects are followed to store the final
resolved URL.

{\def\LTcaptype{none} % do not increment counter
\begin{longtable}[]{@{}ll@{}}
\toprule\noalign{}
Column & Notes \\
\midrule\noalign{}
\endhead
\bottomrule\noalign{}
\endlastfoot
\texttt{video\_id} & FK → \texttt{videos.id} \\
\texttt{url} & Final resolved URL (after following redirects) \\
\texttt{domain} & Extracted domain (e.g. \texttt{bbc.com}) \\
\texttt{title} & Page title fetched at discovery time \\
\texttt{found\_in} & \texttt{description} or \texttt{pinned-comment} \\
\end{longtable}
}

Unique on \texttt{(video\_id,\ url)}.

\phantomsection\label{file-gemini-kg}
\section{Gemini --- Knowledge Graph
Extraction}\label{gemini--knowledge-graph-extraction}

\textbf{File:} \texttt{sentiment\_analysis/articles/extractor.py}
\textbf{Called from:}
\texttt{sentiment\_analysis/articles/service.py\ →\ ingest\_url()}
\textbf{Trigger:} Every time an article is ingested

\begin{center}\rule{0.5\linewidth}{0.5pt}\end{center}

\subsection{What It Does}\label{what-it-does-1}

After an article is scraped and saved, its body text is sent to Google
Gemini to extract a structured knowledge graph. The model identifies
named entities and the relationships between them, which are then stored
in \texttt{articles.kg\_elements} and \texttt{articles.kg\_relations}.

\begin{center}\rule{0.5\linewidth}{0.5pt}\end{center}

\subsection{Models Used}\label{models-used}

Two models are configured --- a primary and a fallback. If the primary
fails for any reason (rate limit, timeout, API error), the request is
automatically retried with the fallback.

{\def\LTcaptype{none} % do not increment counter
\begin{longtable}[]{@{}lll@{}}
\toprule\noalign{}
Role & Model & Config Variable \\
\midrule\noalign{}
\endhead
\bottomrule\noalign{}
\endlastfoot
Primary & \texttt{gemini-3.1-pro-preview} & \texttt{GEMINI\_MODEL} \\
Fallback & \texttt{gemini-2.5-pro} & \texttt{GEMINI\_FALLBACK\_MODEL} \\
\end{longtable}
}

Both are overridable via environment variables.

\begin{center}\rule{0.5\linewidth}{0.5pt}\end{center}

\subsection{What the Model Extracts}\label{what-the-model-extracts}

The model is prompted to return two things:

\textbf{Elements (nodes)} --- named entities with canonical names and
one of six fixed types:

{\def\LTcaptype{none} % do not increment counter
\begin{longtable}[]{@{}ll@{}}
\toprule\noalign{}
Type & Examples \\
\midrule\noalign{}
\endhead
\bottomrule\noalign{}
\endlastfoot
\texttt{PERSON} & Narendra Modi, Hillary Clinton \\
\texttt{ORGANIZATION} & United Nations, NATO \\
\texttt{LOCATION} & New Delhi, the Himalayas \\
\texttt{COUNTRY} & India, Bangladesh \\
\texttt{COMPANY} & Apple, Reliance Industries \\
\texttt{EVENT} & G20 Summit, 2024 General Elections \\
\end{longtable}
}

\textbf{Relations (edges)} --- directed triples connecting two elements:

\begin{verbatim}
{ "subject": "Narendra Modi", "relation": "VISITED", "object": "Bangladesh" }
\end{verbatim}

Relations use standardised uppercase snake\_case labels
(\texttt{APPOINTED\_BY}, \texttt{ACQUIRED}, \texttt{HEADQUARTERED\_IN},
etc.) so they\textquotesingle re reusable across articles.

\begin{center}\rule{0.5\linewidth}{0.5pt}\end{center}

\subsection{How the API is Called}\label{how-the-api-is-called}

\begin{Shaded}
\begin{Highlighting}[]
\NormalTok{response }\OperatorTok{=} \ControlFlowTok{await}\NormalTok{ client.aio.models.generate\_content(}
\NormalTok{    model}\OperatorTok{=}\NormalTok{model,}
\NormalTok{    contents}\OperatorTok{=}\NormalTok{prompt,}
\NormalTok{    config}\OperatorTok{=}\NormalTok{types.GenerateContentConfig(}
\NormalTok{        response\_mime\_type}\OperatorTok{=}\StringTok{"application/json"}\NormalTok{,}
\NormalTok{        response\_schema}\OperatorTok{=}\NormalTok{\_RESPONSE\_SCHEMA,}
\NormalTok{        temperature}\OperatorTok{=}\FloatTok{0.0}\NormalTok{,}
\NormalTok{    ),}
\NormalTok{)}
\end{Highlighting}
\end{Shaded}

\begin{itemize}
\tightlist
\item
  \texttt{response\_mime\_type="application/json"} --- forces structured
  JSON output
\item
  \texttt{response\_schema} --- enforces the exact shape of the response
  including the entity type enum, so the model can\textquotesingle t
  return unexpected values
\item
  \texttt{temperature=0.0} --- deterministic output, minimises
  hallucinations
\end{itemize}

The call is made asynchronously via
\texttt{client.aio.models.generate\_content}.

\begin{center}\rule{0.5\linewidth}{0.5pt}\end{center}

\subsection{Token Truncation}\label{token-truncation}

Before the prompt is built, the body text is truncated to
\texttt{GEMINI\_MAX\_TOKENS} (default: 8000) using \textbf{tiktoken}
with the \texttt{cl100k\_base} encoding. This ensures the request stays
within model limits for long articles. A warning is logged when
truncation occurs.

\begin{center}\rule{0.5\linewidth}{0.5pt}\end{center}

\subsection{Prompt Design}\label{prompt-design}

The prompt instructs the model to think in terms of a formal graph
(nodes and edges). Key rules enforced in the prompt:

\begin{itemize}
\tightlist
\item
  Use canonical, globally recognisable names --- no pronouns or
  abbreviations
\item
  Deduplicate coreferences (same entity mentioned multiple ways = one
  node)
\item
  Only create relations between entities explicitly extracted as
  elements
\item
  Avoid orphaned nodes with no edges unless the entity is the central
  subject
\item
  Return raw JSON only --- no markdown, no explanation
\end{itemize}

\phantomsection\label{file-marianmt-translation}
\section{MarianMT --- Bangla
Translation}\label{marianmt--bangla-translation}

\textbf{File:} \texttt{sentiment\_analysis/articles/translator.py}
\textbf{Called from:}
\texttt{sentiment\_analysis/articles/service.py\ →\ ingest\_url()}
\textbf{Trigger:} Only when \texttt{langdetect} identifies the article
language as Bangla (\texttt{bn})

\begin{center}\rule{0.5\linewidth}{0.5pt}\end{center}

\subsection{What It Does}\label{what-it-does-2}

Translates Bangla-language article content into English before it is
saved to the database and sent to Gemini. The original Bangla text is
preserved alongside the translation in \texttt{*\_original} columns.

Three fields are translated:

\begin{itemize}
\tightlist
\item
  \texttt{headline}
\item
  \texttt{body\_text}
\item
  \texttt{reporters} (each name individually)
\end{itemize}

\begin{center}\rule{0.5\linewidth}{0.5pt}\end{center}

\subsection{Model Used}\label{model-used}

{\def\LTcaptype{none} % do not increment counter
\begin{longtable}[]{@{}ll@{}}
\toprule\noalign{}
Detail & Value \\
\midrule\noalign{}
\endhead
\bottomrule\noalign{}
\endlastfoot
Model & \texttt{Helsinki-NLP/opus-mt-bn-en} \\
Type & MarianMT (sequence-to-sequence transformer) \\
Source & HuggingFace Transformers \\
Config variable & \texttt{TRANSLATION\_MODEL} \\
\end{longtable}
}

The model is loaded \textbf{lazily} --- it is not imported or
initialised at startup, only on the first Bangla article encountered.
This avoids loading several hundred MB of weights on every run.

\begin{center}\rule{0.5\linewidth}{0.5pt}\end{center}

\subsection{Token Limit Handling}\label{token-limit-handling}

MarianMT has a hard input limit of \textbf{512 tokens}. Bangla articles
routinely exceed this, so text is split into chunks before translating:

\begin{enumerate}
\def\labelenumi{\arabic{enumi}.}
\tightlist
\item
  Split on double newlines (paragraph boundaries)
\item
  If a paragraph exceeds 1000 characters (\textasciitilde400--500 Bangla
  tokens), split further on sentence-ending punctuation (\texttt{।},
  \texttt{?}, \texttt{!})
\item
  Accumulate sentences into chunks that stay under 1000 characters
\item
  Each chunk is translated independently and the results are rejoined
  with double newlines
\end{enumerate}

\begin{center}\rule{0.5\linewidth}{0.5pt}\end{center}

\subsection{How the Model is Called}\label{how-the-model-is-called}

\begin{Shaded}
\begin{Highlighting}[]
\NormalTok{inputs }\OperatorTok{=}\NormalTok{ \_tokenizer(}
\NormalTok{    [chunk],}
\NormalTok{    return\_tensors}\OperatorTok{=}\StringTok{"pt"}\NormalTok{,}
\NormalTok{    padding}\OperatorTok{=}\VariableTok{True}\NormalTok{,}
\NormalTok{    truncation}\OperatorTok{=}\VariableTok{True}\NormalTok{,}
\NormalTok{    max\_length}\OperatorTok{=}\DecValTok{512}\NormalTok{,}
\NormalTok{)}
\NormalTok{outputs }\OperatorTok{=}\NormalTok{ \_model.generate(}\OperatorTok{**}\NormalTok{inputs)}
\NormalTok{translated }\OperatorTok{=}\NormalTok{ \_tokenizer.decode(outputs[}\DecValTok{0}\NormalTok{], skip\_special\_tokens}\OperatorTok{=}\VariableTok{True}\NormalTok{)}
\end{Highlighting}
\end{Shaded}

Each chunk goes through the tokenizer → model → decoder independently.
The final output is the chunks joined back together.

\begin{center}\rule{0.5\linewidth}{0.5pt}\end{center}

\subsection{What Gets Stored}\label{what-gets-stored}

{\def\LTcaptype{none} % do not increment counter
\begin{longtable}[]{@{}ll@{}}
\toprule\noalign{}
DB Column & Content \\
\midrule\noalign{}
\endhead
\bottomrule\noalign{}
\endlastfoot
\texttt{headline} & Translated English headline \\
\texttt{body\_text} & Translated English body \\
\texttt{reporters} & Translated English names \\
\texttt{headline\_original} & Original Bangla headline \\
\texttt{body\_text\_original} & Original Bangla body \\
\texttt{reporters\_original} & Original Bangla reporter names \\
\texttt{language} & Set to \texttt{"bn"} to record the source
language \\
\end{longtable}
}

\section{Claude Context}\label{claude-context}

\subsection{Paths}\label{paths}

\begin{itemize}
\tightlist
\item
  Project code: /Users/amaanirfan/Desktop/Data Scraping/Sentiment
  Analysis
\item
  Documentation vault: /Users/amaanirfan/Desktop/Sentiment Analysis Docs
\end{itemize}

\subsection{My Job}\label{my-job}

When Amaan describes work or asks me to document something:

\begin{enumerate}
\def\labelenumi{\arabic{enumi}.}
\tightlist
\item
  Read the relevant files from the project code path above
\item
  Extract context --- file structure, functions, logic, recent changes
\item
  Write clean documentation into the Obsidian vault
\item
  Use proper frontmatter, {[}{[}wikilinks{]}{]}, and tags on every note
\item
  You have to go through the code base and document everything
  accordingly. It should be well structured and nicely written.
\end{enumerate}

\end{document}
